% !TeX program = lualatex
% !TeX encoding = UTF-8

% SoftwareX / Elsevier submission layout for the Splatone implementation paper.
% The previous jlreq-based manuscript is preserved as main-jlreq.tex.
% To avoid changing the manuscript text during the style migration, this file
% extracts the body of main-jlreq.tex and typesets it with Elsevier's
% elsarticle class.

\documentclass[preprint,12pt]{elsarticle}

% Japanese is retained for the current manuscript body. SoftwareX submission
% itself should ultimately be prepared in English.
\usepackage{luatexja}
\usepackage{luatexja-fontspec}

\usepackage{amsmath,amssymb,amsthm}
\usepackage{graphicx}
\usepackage{tikz}
\usetikzlibrary{arrows.meta,positioning}
\usepackage{url}
\usepackage{hyperref}
\usepackage{listings}
\usepackage{xcolor}
\usepackage{booktabs}
\usepackage{tabularx}
\usepackage{array}

\lstdefinestyle{commandline}{
  basicstyle=\ttfamily\small,
  backgroundcolor=\color{black!5},
  frame=single,
  columns=fullflexible,
  keepspaces=true,
  escapeinside={}
}

\lstnewenvironment{commandline}[1][]{
  \lstset{style=commandline,#1}
}{}

\lstdefinestyle{commandstdout}{
  basicstyle=\ttfamily\small\color{white},
  backgroundcolor=\color{darkgray},
  frame=single,
  columns=fullflexible,
  keepspaces=true,
  escapeinside={(*@}{@*)}
}

\lstnewenvironment{commandstdout}[1][]{
  \lstset{style=commandstdout,#1}
}{}

\journal{SoftwareX}

\begin{document}

\begin{frontmatter}

\title{Splatone: A modular platform for multi-source, category-aware geospatial data collection and visualization}

\author[tmu]{Shohei Yokoyama}
\address[tmu]{Tokyo Metropolitan University, 6-6 Asahigaoka, Hino, Tokyo, Japan}

\begin{abstract}
Geotagged photographs, social media posts, and point-of-interest records are valuable sources for studying urban space, tourism, and human activity, but their volume, spatial heterogeneity, and category overlap make them difficult to inspect with simple point maps. Splatone is an open-source Node.js platform for collecting, normalizing, and visualizing such geospatial data. Its architecture separates data acquisition and visualization into independent Provider and Visualizer plugins, allowing new data sources and analytical views to be added without modifying the core pipeline. Existing providers support Flickr, Google Places, and OpenStreetMap/Overpass, and convert heterogeneous responses into a common GeoJSON-based data model. The platform supports category-aware color assignment and visualizers such as point maps, marker clustering, heatmaps, hexagonal aggregation, pie charts, Voronoi diagrams, and DBSCAN-based cluster boundaries. Splatone enables researchers to compare thematic spatial distributions interactively and export reproducible visualization results for further analysis.
\end{abstract}

\begin{keyword}
geosocial data \sep geotagged images \sep geospatial visualization \sep color visualization \sep spatial clustering \sep open-source software
\end{keyword}

\end{frontmatter}

% ---------------------------------------------------------------------------
% SoftwareX code metadata table (Article Template v4 style)
% ---------------------------------------------------------------------------
\section*{Code metadata}

\renewcommand{\arraystretch}{1.2}
\begin{tabularx}{\textwidth}{>{\raggedright\arraybackslash}p{0.08\textwidth} >{\raggedright\arraybackslash}p{0.38\textwidth} X}
\toprule
Nr. & Code metadata description & Metadata \\
\midrule
C1 & Current code version & 0.0.34 \\
C2 & Permanent link to code/repository used for this code version & \url{https://github.com/YokoyamaLab/Splatone/tree/ac66ae6fba442929315cd02df42bc153eb4b46d4} \\
C3 & Permanent link to reproducible capsule & N/A \\
C4 & Legal code license & MIT License \\
C5 & Code versioning system used & Git \\
C6 & Software code languages, tools, and services used & JavaScript (ES modules), Node.js, Express, EJS, Turf.js and related npm packages \\
C7 & Compilation requirements, operating environments \& dependencies & Node.js and npm; dependencies are specified in \texttt{package.json} and \texttt{package-lock.json} \\
C8 & If available, link to developer documentation/manual & \url{https://github.com/YokoyamaLab/Splatone} \\
C9 & Support email for questions & \textit{TODO: add the support/corresponding-author email before submission} \\
\bottomrule
\end{tabularx}

\vspace{1em}

% ---------------------------------------------------------------------------
% Reuse the manuscript body without rewriting it during this format change.
% The extracted body is kept in a normal input file so listings/verbatim content
% is read with the correct catcodes.
% ---------------------------------------------------------------------------
\input{main-body.tex}

% SoftwareX/Elsevier uses a numbered reference style.
\bibliographystyle{elsarticle-num}
\bibliography{references}

\end{document}



